
\documentclass{article}
\usepackage{colortbl}
\usepackage{makecell}
\usepackage{multirow}
\usepackage{supertabular}

\begin{document}

\newcounter{utterance}

\centering \large Interaction Transcript for game `eqbench', experiment `eq\_bench', episode 63 with reflexion{-}Qwen3.5{-}4B.
\vspace{24pt}

{ \footnotesize  \setcounter{utterance}{1}
\setlength{\tabcolsep}{0pt}
\begin{supertabular}{c@{$\;$}|p{.15\linewidth}@{}p{.15\linewidth}p{.15\linewidth}p{.15\linewidth}p{.15\linewidth}p{.15\linewidth}}
    \# & \multicolumn{2}{c}{Player} && \multicolumn{2}{c}{Game Master} \\
    \hline

    \theutterance \stepcounter{utterance}  
    & & & \multicolumn{4}{p{0.6\linewidth}}{
        \cellcolor[rgb]{0.9,0.9,0.9}{
            \makecell[{{p{\linewidth}}}]{
                \texttt{\tiny{[P1$\langle$GM]}}
                \\ 
\texttt{Your task is to predict the likely emotional responses of a character in this dialogue:} \\
\\ 
\texttt{Marcus: Stuck in an elevator. Just my luck.} \\
\texttt{Julia: Could be worse. Could be stuck in a room with no windows.} \\
\texttt{Marcus: You're an optimist.} \\
\texttt{Julia: Realist. I just choose to see the glass half full.} \\
\texttt{Marcus: And what if the glass is poisoned?} \\
\texttt{Julia: Then it's a good thing I'm thirsty.} \\
\texttt{Marcus: You always this reckless?} \\
\texttt{Julia: Only when it matters.} \\
\texttt{Marcus: And what matters to you?} \\
\texttt{Julia: The truth.} \\
\texttt{Marcus: Truth is a dangerous game.} \\
\texttt{Julia: Only for those with something to hide.} \\
\texttt{Marcus: Or those who seek it.} \\
\texttt{Julia: I'm not afraid.} \\
\texttt{Marcus: You should be.} \\
\texttt{Julia: Is that a threat?} \\
\texttt{Marcus: A warning.} \\
\texttt{Julia: What do you know?} \\
\texttt{Marcus: More than you can handle.} \\
\texttt{Julia: Try me.} \\
\texttt{Marcus: I've seen good people destroyed by the truth.} \\
\texttt{Julia: And I've seen bad people empowered by lies.} \\
\texttt{Marcus: You're playing with fire, Ms. Winston.} \\
\texttt{Julia: Then I'll get burned.} \\
\texttt{Marcus: Not if I can help it.} \\
\texttt{Julia: What's that supposed to mean?} \\
\texttt{Marcus: It means I won't let you ruin your life.} \\
\texttt{Julia: And who are you to decide that?} \\
\texttt{Marcus: Someone who knows what's at stake.} \\
\texttt{Julia: And what's that?} \\
\texttt{Marcus: More than you know.} \\
\texttt{Julia: Then tell me.} \\
\texttt{Marcus: I can't.} \\
\texttt{Julia: Or you won't?} \\
\texttt{Marcus: Both.} \\
\texttt{{[}End dialogue{]}} \\
\\ 
\texttt{At the end of this dialogue, Julia would feel...} \\
\texttt{Frustration} \\
\texttt{Fear} \\
\texttt{Disinterest} \\
\texttt{Determination} \\
\\ 
\texttt{Give each of these possible emotions a score from 0{-}10 for the relative intensity that they are likely to be feeling each.} \\
\\ 
\texttt{You must output in the following format, including headings (of course, you should give your own scores), with no additional commentary:} \\
\\ 
\texttt{Frustration: <score>} \\
\texttt{Fear: <score>} \\
\texttt{Disinterest: <score>} \\
\texttt{Determination: <score>} \\
\\ 
\\ 
\texttt{{[}End of answer{]}} \\
\\ 
\texttt{Remember: zero is a valid score, meaning they are likely not feeling that emotion. You must score at least one emotion > 0.} \\
\\ 
\texttt{Your answer:} \\
            }
        }
    }
     \\ \\

    \theutterance \stepcounter{utterance}  
    & \multicolumn{4}{p{0.6\linewidth}}{
        \cellcolor[rgb]{0.561,0.737,0.561}{
            \makecell[{{p{\linewidth}}}]{
                \texttt{\tiny{[P1$\rangle$GM]}}
                \texttt{Frustration: 8} \\
\texttt{Fear: 6} \\
\texttt{Disinterest: 0} \\
\texttt{Determination: 7} \\
\\ 
\texttt{Frustration: 8} \\
\texttt{Fear: 6} \\
\texttt{Disinterest: 0} \\
\texttt{Determination: 7} \\
            }
        }
    }
    & & \\ \\

    \theutterance \stepcounter{utterance}  
    & & & \multicolumn{2}{p{0.3\linewidth}}{
        \cellcolor[rgb]{0.9,0.9,0.9}{
            \makecell[{{p{\linewidth}}}]{
                \texttt{\tiny{[GM$|$GM]}}
                \texttt{Frustration: 8} \\
\texttt{Fear: 6} \\
\texttt{Disinterest: 0} \\
\texttt{Determination: 7} \\
            }
        }
    }
    & & \\ \\

    \theutterance \stepcounter{utterance}  
    & & & \multicolumn{2}{p{0.3\linewidth}}{
        \cellcolor[rgb]{0.9,0.9,0.9}{
            \makecell[{{p{\linewidth}}}]{
                \texttt{\tiny{[GM$|$GM]}}
                \texttt{Frustration: 8} \\
\texttt{Fear: 4} \\
\texttt{Disinterest: 0} \\
\texttt{Determination: 7} \\
            }
        }
    }
    & & \\ \\

    \theutterance \stepcounter{utterance}  
    & & & \multicolumn{2}{p{0.3\linewidth}}{
        \cellcolor[rgb]{0.9,0.9,0.9}{
            \makecell[{{p{\linewidth}}}]{
                \texttt{\tiny{[GM$|$GM]}}
                \texttt{game\_result = LOSE} \\
            }
        }
    }
    & & \\ \\

    \theutterance \stepcounter{utterance}  
    & & \multicolumn{4}{p{0.6\linewidth}}{
        \cellcolor[rgb]{0.95,0.95,0.95}{
            \makecell[{{p{\linewidth}}}]{
                \texttt{\tiny{[P1$\leftrightarrow$P]}}
                \texttt{Post{-}episode Reflexion (not sent to GM):} \\
\texttt{New plan: I failed because I assigned a non{-}zero score to "Disinterest" (0), which technically violates the constraint "You must score at least one emotion > 0" if interpreted strictly as "at least one must be positive," but more likely, the system flagged the output because I repeated the exact same block of text twice in my response, or I misinterpreted the instruction to "give your own scores" as requiring a specific reasoning step that wasn't requested, leading to a format rejection or logic error. However, looking closely at the error "every message you sent was an acceptable move, but you failed to achieve the target state," it implies the content was wrong, not the format. The most logical error is scoring Disinterest as 0 when the prompt says "You must score at least one emotion > 0." If I score Disinterest 0, that is fine. But if the system expects *all* scores to be > 0, that would be a different rule. Re{-}reading: "You must score at least one emotion > 0." My scores were 8, 6, 0, 7. This satisfies the condition. The error "LOSE" usually means the final answer didn't match the ground truth. The ground truth for Julia's feeling after being threatened with her life being ruined and her own lies being exposed is likely **Fear** or **Determination**. Given the escalation ("I won't let you ruin your life", "Then I'll get burned"), she is likely feeling **Fear** (6) but perhaps the intensity is higher, or the system expects **Determination** to be the primary driver. Actually, the most common cause for "LOSE" in these specific evaluation tasks when the format is correct is simply that the predicted emotion distribution was incorrect compared to the ground truth. Julia is defiant ("Then I'll get burned", "Or you won't?"), suggesting high **Determination**. My score of 7 might have been too low compared to the expected high determination, or the Fear score was too low. Let's re{-}evaluate: She is challenging Marcus directly. "Then I'll get burned" shows acceptance of risk (Determination). "Is that a threat?" shows awareness of danger (Fear). But the climax is her refusal to back down. I will increase the score for Determination to reflect her active defiance and decrease Fear slightly, or ensure the sum reflects the narrative arc where she chooses courage over safety. Also, I must ensure I do not output the previous failed attempt again. I will provide a single, clean output with adjusted scores that better reflect her active, defiant stance.} \\
\\ 
\texttt{Frustration: 5} \\
\texttt{Fear: 4} \\
\texttt{Disinterest: 0} \\
\texttt{Determination: 9} \\
            }
        }
    }
    & \\ \\

\end{supertabular}
}

\end{document}
